\begin{helper}
Let \(A,t,q\) be nonnegative integers with \(t\ge1\) and \(q\ge1\).  Define
\[
  w=\left\lfloor A q\log q\right\rfloor .
\]
Assume that the truncated natural-number exponent satisfies
\[
  32tq\log(4q^t)<\frac{w}{t+1}-1 .
\]
Then
\[
  w\le A q\log q
\]
and
\[
  \left(1-\frac{1}{32tq}\right)^{w/(t+1)-1}
  <\frac{1}{4q^t}.
\]
Here \(w/(t+1)-1\) denotes natural-number subtraction after integer division,
as in the Lean statement.
\end{helper}
\begin{proof}
Since \(q\ge1\), we have \(\log q\ge0\).  Thus
\[
  A q\log q\ge0.
\]
For a nonnegative real number \(x\), its natural-number floor satisfies
\(\lfloor x\rfloor\le x\).  Applying this to \(x=Aq\log q\) gives
\[
  w=\lfloor Aq\log q\rfloor\le Aq\log q.
\]

Put
\[
  e=\frac{w}{t+1}-1,\qquad D=32tq,\qquad M=4q^t.
\]
The hypotheses \(t\ge1\) and \(q\ge1\) give \(D>0\), \(D\ge1\), and
\(M>0\).  Hence \(0\le 1/D\le1\), so \(1-1/D\ge0\).  The elementary
exponential inequality \(1-x\le e^{-x}\), applied with \(x=1/D\), gives
\[
  1-\frac1D\le \exp(-1/D).
\]
Raising both sides to the nonnegative integer power \(e\) gives
\[
  \left(1-\frac1D\right)^e\le \exp(-1/D)^e
  =\exp(-e/D).
\]

The assumed lower bound \(D\log M<e\), divided by the positive number \(D\),
gives
\[
  \log M<e/D.
\]
Equivalently,
\[
  -e/D<-\log M.
\]
Since the exponential function is strictly increasing and \(M>0\),
\[
  \exp(-e/D)<\exp(-\log M)=\frac1M.
\]
Combining the last two displayed inequalities yields
\[
  \left(1-\frac1D\right)^e<\frac1M.
\]
Substituting back \(D=32tq\), \(M=4q^t\), and
\(e=w/(t+1)-1\) gives the required shrink threshold.
\end{proof}
